\section{Discussion}
\label{sec:discussion}

\subsection{Practical Applications}
\label{sec:practical_applications}

We present \system{} as a practical framework that can be built on commodity devices such as smartphones, smartwatches, and earbuds. For example, \system{} can support context-aware assistance by integrating IMU from smartphones and smartwatches to infer situations (e.g., driving, being in the office, staying at home), and health monitoring by jointly interpreting ECG and PPG to provide early feedback on abnormal conditions. In these scenarios, synchronized sensor streams can be collected, segmented into recent sliding windows, and converted into standardized features using established preprocessing pipelines~\citep{wesad}. To facilitate practical adoption, Appendix~\ref{sec:appendix_prompts} provides task- and modality-agnostic prompt configurations, allowing practitioners to adapt \system{} to new applications by replacing placeholders with their own task definitions, modality descriptions, and sensor metadata.

\subsection{Uncertainty-Aware Fusion}

\system{} addresses sensor uncertainty implicitly through semantic interpretation and statistical agreement. This enables arbitration without requiring additional supervision or external reliability models. At the same time, explicit reliability-aware fusion could further improve performance by incorporating reliability signals as fusion weights. Promising directions include confidence-weighted voting, tool-augmented modality agents that estimate signal quality through data processing tools, and historical reliability modeling when data streams are available. We leave these extensions to future work and focus here on establishing a general, training- and model-free hybrid reasoning protocol under unknown uncertainty.